% from main.tex
%%%%%%%%%%  scan idx 151  |  folio 138  %%%%%%%%%%
Cette formule est le point de départ de la théorie des intersections sur les schémas. On notera que par contre le produit tensoriel (dérivé) de deux complexes pseudo-cohérents à cohomologie bornée n'est pas en général à cohomologie bornée.

\textbf{Corollaire 5.9.} Supposons que \(S=(\mathrm{Ens})\) (donc \(A\) «~est~» un anneau ordinaire) et que tout \(A\)-module à gauche soit de tor-dimension finie (c'est le cas par exemple si \(A\) est de dimension cohomologique globale finie, voir (MULT VI 2.6) et (EGA \(0_{\mathrm{IV}}\) 17.2.8)). Alors on a
\[
\begin{aligned}
\text{(i)}\quad & \Dpx^b({}_A S)=\Dpx({}_A S)_{\torf},\quad\text{et}\\
\text{(ii)}\quad & \Dpx^b({}_A S)_{\mathrm{coh}}=\Dpx({}_A S)_{\mathrm{parf}}.
\end{aligned}
\]
Autrement dit, pour que \(E\in\Ob\Dpx({}_A S)\) soit de tor-dimension finie (resp. de tor-dimension finie et parfait) il faut et il suffit que \(E\) soit à cohomologie bornée (resp. à cohomologie bornée et pseudo-cohérent).

\emph{Preuve.} D'après (5.8), il suffit de prouver (i). Soit \(E\in\Ob\Dpx^b({}_A S)\). On peut supposer (quitte à prolonger \(E\) par \(0\)) que \(E\in\Ob\Dpx^b({}_A S)\), où \(S\) est l'objet final de \(S\). Procédant par récurrence sur le nombre d'objets de cohomologie non nuls de \(E\), on se ramène, par dévissage, au cas où \(E\) est un \(A\)-Module ordinaire, et l'on gagne.

On peut énoncer plus généralement :

\textbf{Corollaire 5.10.} Soit \(X\in\Ob S\). Supposons que \(S_{/X}\) ait assez de points (SGA 4 IV 6.4.1.) et que, pour tout point \(s\) de \(S_{/X}\), tout \(A_s\)-module à gauche soit de tor-dimension finie. Alors on a
\[
  \Dpx^b({}_A X)_{\mathrm{coh}}=\Dpx^b({}_A X)_{\mathrm{parf}}.
\]
